\chapter{Les notions à connaître}
\label{chap:notions}

\begin{chapeau}
Ce chapitre donne le vocabulaire utilisé partout ensuite : l'anatomie du membre pelvien, la façon de décrire une
déformité, l'ostéotomie et ses outils, puis le scanner et le format DICOM. Un lecteur qui connaît déjà la chirurgie
orthopédique vétérinaire peut passer directement à la \cref{sec:scanner}.
\end{chapeau}

\section{Le membre pelvien du chien}

Le membre pelvien (la patte arrière) comprend, de haut en bas :
\begin{itemize}
  \item le \textbf{bassin}, dont la cavité (l'\emph{acétabulum}) reçoit la tête du fémur : c'est la hanche ;
  \item le \textbf{fémur}, l'os de la cuisse. En haut, sa \emph{tête} sphérique, son \emph{col} et le \emph{grand
  trochanter} ; en bas, deux \emph{condyles} séparés par l'\emph{échancrure intercondylaire}, et devant eux la
  \emph{trochlée}, la gouttière où glisse la rotule ;
  \item la \textbf{rotule} (\emph{patella}), petit os pris dans le tendon du quadriceps, et deux \textbf{fabelles},
  petits os placés derrière les condyles, dans les tendons du muscle gastrocnémien ;
  \item le \textbf{tibia}, l'os de la jambe. En haut, le \emph{plateau tibial} (deux surfaces articulaires séparées
  par l'\emph{éminence intercondylaire}) et, devant, la \emph{tubérosité tibiale}, où s'attache le tendon de la
  rotule ; en bas, la \emph{cochlée}, surface articulaire de la cheville, et la \emph{malléole médiale} ;
  \item la \textbf{fibula}, os long et fin, le long du tibia, du côté latéral et un peu en arrière ; son extrémité
  basse forme la \emph{malléole latérale} ;
  \item le \textbf{tarse} (dont le \emph{talus} et le \emph{calcanéus}), les \textbf{métatarses} et les
  \textbf{phalanges}.
\end{itemize}

Un os long a trois zones. La \emph{diaphyse} est le corps de l'os : un tube d'os dense, la \emph{corticale}, autour du
\emph{canal médullaire}. Les \emph{épiphyses} sont les extrémités, faites d'os spongieux sous une fine corticale et
recouvertes de \emph{cartilage articulaire}. Entre les deux, la \emph{métaphyse} s'élargit. Chez le jeune chien, la
\emph{physe} (cartilage de croissance) sépare épiphyse et métaphyse ; une coupe à travers une physe ouverte arrête la
croissance.

\section{Les directions et les plans}
\label{sec:directions}

En anatomie vétérinaire, on ne dit pas « devant » et « derrière », mais :
\begin{itemize}
  \item \textbf{crânial} (vers la tête, l'avant du membre) et \textbf{caudal} (vers la queue, l'arrière) ;
  \item \textbf{proximal} (vers le tronc, le haut de l'os) et \textbf{distal} (vers le pied, le bas de l'os) ;
  \item \textbf{médial} (vers le plan de symétrie du corps, l'intérieur) et \textbf{latéral} (vers l'extérieur).
\end{itemize}

Trois plans servent à regarder un os :
\begin{itemize}
  \item le \textbf{plan frontal} : on regarde l'os de face (ou de dos), comme sur une radiographie
  crânio-caudale ;
  \item le \textbf{plan sagittal} : on regarde l'os de profil, comme sur une radiographie médio-latérale ;
  \item le \textbf{plan transverse} : on regarde une coupe perpendiculaire à l'os, comme une coupe de scanner.
\end{itemize}

\begin{aretenir}
Dans tout le logiciel, chaque os est exprimé dans \textbf{son propre repère} : $X$ latéral, $Y$ proximal, $Z$ caudal,
en millimètres. Le plan frontal est donc le plan $(X, Y)$, le plan sagittal le plan $(Z, Y)$. Un membre gauche est
transformé en son image dans un miroir, pour être traité exactement comme un membre droit
(\cref{sec:repere-os}).
\end{aretenir}

\section{Décrire une déformité}
\label{sec:deformite}

\subsection{Les trois composantes}

Une déformité angulaire se décrit dans les trois plans :
\begin{itemize}
  \item \textbf{de face} : le bas de l'os part vers l'intérieur (\emph{varus}) ou vers l'extérieur (\emph{valgus}) ;
  \item \textbf{de profil} : le sommet de la courbure est en avant (\emph{procurvatum}) ou en arrière
  (\emph{recurvatum}) ;
  \item \textbf{en torsion} : le bas de l'os est tourné autour de l'axe de l'os par rapport au haut
  (\emph{torsion interne} ou \emph{externe}).
\end{itemize}

Le point où l'os change de direction s'appelle l'\emph{apex}. Une déformité à un seul apex est dite
\emph{uniapicale} ; un os tordu à deux endroits est \emph{biapical}. Ce projet traite les déformités uniapicales.

\subsection{Lignes articulaires et axes}

Pour mesurer une déformité, on trace deux sortes de lignes :
\begin{itemize}
  \item une \textbf{ligne d'orientation articulaire} (JOL, pour \emph{joint orientation line}) passe par deux repères
  d'une articulation, par exemple les points les plus bas des deux condyles du fémur ;
  \item l'\textbf{axe} d'un os ou d'un segment d'os. L'\emph{axe mécanique} relie les centres des deux articulations ;
  l'\emph{axe anatomique} suit le milieu de la diaphyse. Au tibia du chien, on utilise l'axe mécanique ; au fémur,
  l'axe anatomique \cite{tomlinson2007,dismukes2007}.
\end{itemize}

L'angle entre une ligne articulaire et l'axe est un \textbf{angle articulaire}. Son nom se lit lettre par lettre :
\textbf{m} ou \textbf{a} (axe mécanique ou anatomique), \textbf{L} ou \textbf{M} (angle mesuré du côté latéral ou
médial), \textbf{P} ou \textbf{D} (articulation proximale ou distale), \textbf{F} ou \textbf{T} (fémur ou tibia),
\textbf{A} (angle). Dans le plan sagittal, \textbf{Cr} et \textbf{Ca} (ou \textbf{Cd}) indiquent crânial et caudal.

\begin{table}[htbp]
\centering\small
\caption{Les angles articulaires utilisés par le projet.}
\label{tab:angles}
\begin{tabularx}{\linewidth}{lL{1.6cm}Y}
\toprule
\titrecol{Angle} & \titrecol{Vue} & \titrecol{Définition} \\
\midrule
aLPFA & fémur, face & entre l'axe anatomique et la ligne centre de la tête -- sommet du grand trochanter (côté latéral) \\
aLDFA & fémur, face & entre l'axe anatomique et la ligne des points distaux des condyles (côté latéral) \\
aCdDFA & fémur, profil & entre l'axe anatomique et la ligne fond de la trochlée -- milieu des points caudaux des condyles \\
mMPTA & tibia, face & entre l'axe mécanique et la ligne des centres des plateaux (côté médial) \\
mMDTA & tibia, face & entre l'axe mécanique et la ligne des sillons de la cochlée (côté médial) \\
mCaPTA & tibia, profil & entre l'axe mécanique et la ligne bords crânial -- caudal du plateau médial (côté caudal) \\
mCrDTA & tibia, profil & entre l'axe mécanique et la ligne bords crânial -- caudal de la surface distale (côté crânial) \\
\bottomrule
\end{tabularx}
\end{table}

Chez un chien sain, ces angles ont des valeurs connues, publiées pour plusieurs races (\cref{sec:normales}). Un angle
qui s'écarte de plus de deux écarts-types de la valeur normale signale une déformité. La \cref{fig:angles-reels}
montre ces angles mesurés par le logiciel sur les os du lévrier.

\section{L'ostéotomie correctrice}
\label{sec:osteotomie}

\subsection{Couper et redresser}

Pour redresser l'os, le chirurgien coupe l'os à la scie, puis il fait tourner le morceau du bas (le fragment distal)
jusqu'à ce que l'os soit droit. Deux façons de couper sont courantes (\cref{fig:coins}) :
\begin{itemize}
  \item le \textbf{coin de fermeture} : deux traits de scie retirent un petit coin d'os du côté bombé (convexe) ; on
  referme l'os. Les deux surfaces d'os restent en contact, ce qui est stable, mais l'os raccourcit un peu ;
  \item le \textbf{coin d'ouverture} : un seul trait de scie, du côté creux (concave) ; on ouvre l'os. Il s'allonge un
  peu, mais il reste un espace vide qui doit se combler (parfois avec une greffe).
\end{itemize}

\fig[0.44]{coins_fermeture_ouverture.png}{Coin de fermeture et coin d'ouverture, avant et après la correction. La
scie entre toujours du côté opposé à la charnière : c'est là que se pose le guide.}{fig:coins}

Dans les deux cas, on peut garder un petit pont d'os intact de l'autre côté : la \textbf{charnière}. Les deux
fragments tournent autour d'elle, comme autour d'une charnière de porte. La charnière apporte de la stabilité et
évite que les fragments ne glissent l'un sur l'autre. Une coupe qui traverse tout l'os est une \emph{ostéotomie
complète} : il faut alors tenir les fragments autrement, par exemple avec un guide de réduction.

\subsection{Les outils}

\begin{itemize}
  \item La \textbf{scie oscillante} a une lame fine (environ 0,6~mm) ; chaque trait enlève l'épaisseur de la lame.
  \item Les \textbf{broches de Kirschner} sont de fines tiges métalliques (ici 1,6~mm de diamètre) forées dans l'os.
  Elles fixent un guide, matérialisent une direction ou tiennent un fragment.
  \item La \textbf{plaque vissée} fixe définitivement les fragments ; chaque fragment doit recevoir au moins trois vis
  qui traversent les deux corticales.
\end{itemize}

\subsection{Les voies d'abord}

Pour atteindre l'os, le chirurgien ouvre la peau et écarte les muscles par une \emph{voie d'abord}. Toutes les faces
de l'os ne sont pas accessibles :
\begin{itemize}
  \item au \textbf{fémur}, l'abord latéral est le plus courant ; l'abord médial sert au fémur distal. Le côté crânial
  est occupé par le quadriceps et la rotule, le côté caudal par le nerf sciatique et les vaisseaux ;
  \item au \textbf{tibia}, la face médiale est sous la peau, donc facile d'accès ; la face latérale est couverte par
  les muscles crâniaux, et la fibula est derrière.
\end{itemize}
Le logiciel tient compte de ces contraintes pour décider où poser le guide (\cref{sec:abord}).

\section{Le scanner et les images DICOM}
\label{sec:scanner}

\subsection{Le principe du scanner}

Le scanner fait tourner un tube à rayons X autour de l'animal et mesure, sous des centaines d'angles, combien chaque
rayon est atténué. Un calcul de reconstruction en déduit, pour chaque petit cube de l'espace (un \emph{voxel}), un
coefficient d'atténuation $\mu$. On l'exprime en \textbf{unités Hounsfield} \cite{hounsfield1973} :
\begin{equation}
  \mathrm{HU} = 1000 \times \frac{\mu - \mu_{\text{eau}}}{\mu_{\text{eau}} - \mu_{\text{air}}}.
  \label{eq:hu}
\end{equation}
L'eau vaut donc 0~HU et l'air $-1000$~HU. Plus un tissu est dense, plus sa valeur est grande (\cref{tab:hu}).

\begin{table}[htbp]
\centering\small
\caption{Valeurs typiques au scanner (celles utilisées dans \fichier{parametres.py}).}
\label{tab:hu}
\begin{tabular}{lr}
\toprule
\titrecol{Tissu} & \titrecol{HU} \\
\midrule
air & $-1000$ \\
graisse & environ $-100$ \\
eau & 0 \\
tissus mous (muscles) & 0 à 80 \\
os spongieux & 250 à 700 \\
os cortical & 1200 à 1500 \\
\bottomrule
\end{tabular}
\end{table}

Trois défauts de l'image comptent pour la suite :
\begin{itemize}
  \item le \textbf{bruit} : deux voxels du même tissu n'ont pas exactement la même valeur. Dans les tissus mous du
  lévrier, l'écart-type mesuré est d'environ 49~HU (noyau de reconstruction osseux « BONEPLUS ») ;
  \item le \textbf{flou} : le scanner ne voit pas un bord net. Un bord d'os apparaît comme une transition progressive,
  sur environ un millimètre. Ce flou suit à peu près une courbe en cloche (gaussienne) ;
  \item l'\textbf{effet de volume partiel} : un voxel à cheval sur l'os et le tissu voisin prend une valeur
  intermédiaire. Dans une articulation serrée, quelques voxels de valeur moyenne relient deux os.
\end{itemize}

\subsection{Le format DICOM}

Un scanner est livré au format \textbf{DICOM} \cite{dicom}. Chaque coupe est un fichier qui contient l'image et un
en-tête de métadonnées. Le logiciel utilise surtout :
\begin{itemize}
  \item \code{SeriesInstanceUID} : l'identifiant de la série (un examen peut contenir plusieurs séries) ;
  \item \code{ImagePositionPatient} (IPP) : la position, en millimètres, du centre du premier pixel de la coupe ;
  \item \code{ImageOrientationPatient} (IOP) : les directions des lignes et des colonnes de l'image ;
  \item \code{PixelSpacing} : l'écart entre deux lignes et entre deux colonnes (mm) ;
  \item \code{SliceThickness} : l'épaisseur de la coupe ;
  \item \code{RescaleSlope} et \code{RescaleIntercept} : la règle qui transforme les valeurs stockées en HU.
\end{itemize}

Les positions sont données dans le \textbf{repère LPS} du patient : $x$ vers la gauche du patient
(\emph{Left}), $y$ vers l'arrière (\emph{Posterior}), $z$ vers la tête (\emph{Superior}). C'est aussi le repère des
fichiers produits par 3D Slicer \cite{fedorov2012}. Les scanners du lévrier comptent de 442 à 748 coupes de 0,625~mm.

\subsection{Topogramme ou scanner : ce qu'il faut demander}
\label{sec:topogramme}

Avant l'examen, le scanner fait une image de repérage, le \textbf{topogramme} (ou \emph{scout}) : une seule image 2D,
comme une radiographie, qui sert à choisir la zone à scanner. Elle ne suffit pas au logiciel, qui a besoin du
\textbf{volume}, c'est-à-dire de toutes les coupes.

\begin{attention}[title={Un cas vécu}]
Le premier dossier clinique reçu (27/10/2026, déformité angulaire chez un chien de 12 mois) ne contenait qu'une
image : le topogramme latéral, 512 × 512 pixels de 2~mm. À l'export, seule la première série avait été envoyée.
Aucune planification n'est possible avec une seule projection : il faut redemander l'étude complète.
\end{attention}

\begin{aretenir}[title={Ce qu'il faut demander au vétérinaire}]
\begin{itemize}
  \item l'\textbf{étude complète} en DICOM, avec toutes ses séries (pas seulement la série sélectionnée à l'écran) ;
  \item les \textbf{coupes axiales en reconstruction osseuse}, d'épaisseur 1~mm ou moins ;
  \item si possible les \textbf{deux membres} : le membre sain sert de référence pour la correction ;
  \item des fichiers \textbf{rendus anonymes} (sans nom de l'animal ni du propriétaire), et un dossier nommé par un code.
\end{itemize}
\end{aretenir}
